% RESEARCH_DRAFT_SHA256: 04f2dc40c42134e129a41c61d9093af49ecc8603e7c1e71a5bf6121e22ab725c
% Layout and source-reference wrappers expanded for the web renderer.
\documentclass[11pt]{article}
\usepackage[margin=25mm]{geometry}
\usepackage{fontspec,unicode-math}
\setmainfont{Latin Modern Roman}
\setmathfont{Latin Modern Math}
\usepackage{amsmath,amssymb,mathtools,xurl,hyperref}
\hypersetup{colorlinks=true,urlcolor=blue,linkcolor=blue}
\setcounter{secnumdepth}{0}
\setlength{\parindent}{0pt}
\setlength{\parskip}{5pt}
\setlength{\emergencystretch}{3em}
\newcommand{\sref}[2]{\hyperlink{source:#1:#2}{[S#2]}}
\newcommand{\cataloguescope}[1]{\paragraph{Source record.}#1}
\newcommand{\floor}[1]{\left\lfloor#1\right\rfloor}
\begin{document}
\section*{Research-attempt provenance}
This standalone research attempt retains the complete original definition body below,
including its historical source and preparation notes. Those notes describe the
earlier catalogue preparation; the new attempt and its verification are identified
separately after that body. No original ID, catalogue status, or source statement is changed.

\section{Explicit near-minimal units from the golden-ratio radical}\label{statement}

\subsection{Definitions and mathematical statement}

\paragraph{Shared notation.}Standard number systems are $\mathbb N,\mathbb Z,\mathbb Q,\mathbb R,\mathbb C$. Specialized notation and index ranges follow the cited source and the definitions below.


Let $K=\mathbb Q[X]/(X^{312}-X^{156}-1)$ be the source's quotient algebra and let $\mathcal O_K$ denote its elements integral over $\mathbb Z$ (componentwise if $K$ splits into number fields). A unit $u\in\mathcal O_K$ has an inverse in $\mathcal O_K$. Archimedean embeddings give sizes such as heights or norms of logarithmic absolute-value vectors, but different choices produce different meanings of short.
\paragraph{Record scope and qualifications.}
The catalogue title is retained as metadata. Named certificates, contracts, bounds or frameworks in that title are not asserted as results of the original source extract.
The excerpt fixes a polynomial but no height, embedding norm, equivalence modulo roots of unity, input model, or output format. Those choices are required for a quantitative short-unit bound; none is silently selected. The supplied catalogue marks this record partially solved; its present research status was not independently checked.
\paragraph{Statement recorded in the supplied source.}
Compute short units in the ring of integers of $\mathbb Q[x]/(x^{312} - x^{156} -1)$. How short can they be, and how fast can you find them? \href{http://aimpl.org/quantumalg/6/}{[S2]}
\paragraph{Source status.}
The supplied catalogue labels this source record partially solved; that metadata was checked against the supplied snapshot on 2026-10-11. The label is preserved for provenance and is not an independent verification of present open status or mathematical correctness.
\subsection{Short English statement}
The source asks for small algebraic units and efficient algorithms to find them.
\subsection{Sources}
\begin{itemize}
\item[S1] Ulam AI and the UnsolvedMath Contributors, UnsolvedMath: A Curated Collection of Open Mathematics Problems, catalogue entry. Catalogue curation and metadata: CC BY 4.0. The source extract is retained; definitions, English summaries and qualifications are editorial additions.
\url{https://huggingface.co/datasets/ulamai/UnsolvedMath}.
\item[S2] AIM Workshop Problem Lists, source linked by the catalogue.
\url{http://aimpl.org/quantumalg/6/}.
\end{itemize}
\paragraph{Preparation and review.}
The original source and UnsolvedMath contributors retain attribution for the source extract. Codex (GPT-6; exact serving revision not exposed) assisted with definitions, English summaries, scope notes and formatting. No new human mathematical authorship or independently verified solution is asserted. Catalogue identity checked against the complete supplied registry and current repository. Distinct upstream records, including repeated questions and broad or damaged extracts, are retained for complete coverage of the source CS category. A repeated record is not a new independent conjecture. Complete mathematical-equivalence and current-literature review has not been performed. Preparation and final review follow contribution prompts C and D, with source gaps disclosed. Structural and rendering checks do not establish a new theorem or open status.
\label{end:statement}

\section{Literature and context check}
Fresh retrieval of the original AIM quantum-algorithm problem pages failed (web timeout and direct network/DNS failure; an escalated retry of the related chapter-four page also failed certificate validation). The supplied original extract remains the source text; successful fresh access to those pages is not claimed. The original question fixes a quotient polynomial but leaves ``short'' undefined. The following exact computation gives one explicit unit and reports both maximum absolute value and an unweighted Euclidean norm of its full logarithmic embedding vector. These chosen quantities are auxiliary; no claim of a shortest unit in the full ring of integers is made.
\section{Research attempt: An explicit unit and its complete archimedean size vector}
Let $F(X)=X^{312}-X^{156}-1$ and let $u$ be the class of $X$ in $K=\mathbb Q[X]/(F)$. Since $u$ satisfies a monic integer polynomial, it is integral. The defining relation gives
\[
 u\bigl(u^{311}-u^{155}\bigr)=u^{312}-u^{156}=1.
\]
Thus its inverse is an integer polynomial in the integral element $u$, hence integral as well: sums and products of algebraic integers are integral, since the ring generated by finitely many integral elements is a finitely generated abelian module, and multiplication on it yields a monic annihilating polynomial by the determinant identity. This verifies $u\in\mathcal O_K^\times$ without assuming $F$ irreducible.

There are 312 distinct complex roots. Indeed $F'(X)=156X^{155}(2X^{156}-1)$ has no common root with $F$: zero is not a root of $F$, and $X^{156}=1/2$ would give $F=-5/4$. Put $\varphi=(1+\sqrt5)/2$. The two roots of $Y^2-Y-1$ are $\varphi$ and $-1/\varphi$, so exactly 156 roots of $F$ have modulus $\varphi^{1/156}$ and 156 have modulus $\varphi^{-1/156}$. The quotient algebra is a product of number fields because a square-free polynomial factors into distinct irreducibles and the Chinese remainder identity splits the factors. Its complex embeddings, counted over all components, are exactly evaluations at these 312 roots.

Consequently the full logarithmic vector $\lambda(u)$ has 156 entries $\log\varphi/156$ and 156 entries $-\log\varphi/156$. In the declared unweighted norms,
\[
 \|\lambda(u)\|_\infty={\log\varphi\over156},\qquad
 \|\lambda(u)\|_2={\sqrt{312}\log\varphi\over156},\qquad
 \max_\sigma|\sigma(u)|=\varphi^{1/156}.
\]
The sum of log absolute values is zero, consistent with a unit of absolute norm one; the determinant norm itself is $(-1)^{312}F(0)=-1$. The element is not a root of unity because some embedding has absolute value greater than one.

For all integer $k$, $\lambda(u^k)=k\lambda(u)$, so within the explicitly constructed cyclic subgroup $\langle u\rangle$ the shortest nonzero logarithmic vectors are those of $u$ and $u^{-1}$. This statement supplies no lower bound for units outside that subgroup. It also shows why excluding trivial units matters: $1$ has logarithmic length zero.

The same construction for $X^{2m}-X^m-1$ gives an integral unit with inverse $u^{2m-1}-u^{m-1}$ and unweighted log norm $\sqrt{2/m}\log\varphi$. The proof above repeats with $156$ replaced by $m$. This family explains the numerical smallness through its growing number of embeddings. Its sparse polynomial description has only two terms in the inverse, but explicitly listing all embeddings or a dense degree-$2m$ representation has a different output cost. A bound on construction time must say which representation is requested.
\section{Verification and adversarial audit}
The inverse was checked by multiplication in the exact quotient. Separability was proved before using all root embeddings, so reducibility is harmless. Norm conventions count every complex embedding individually; grouping conjugate embeddings with different weights would change the displayed Euclidean norm. The claim of minimality is restricted to powers of the one constructed unit.
\paragraph{Audit corrections.} A small log vector was not promoted to a global shortest-unit theorem. The quotient is treated componentwise when necessary, and the trivial unit is excluded only in the declared auxiliary comparison.
\paragraph{External computation.} Python 3.14.2 evaluated the exact expressions for illustration: $\|\lambda(u)\|_\infty\approx0.00308469118628$, $\|\lambda(u)\|_2\approx0.0544865098074$, and the maximum embedding modulus is approximately $1.003089453742$. The algebraic inverse and exact expressions, not floating-point values, establish the claims.
\section{Scope and remaining gaps}
No fundamental unit system, irreducibility theorem, global shortest-unit result, or fully specified bit-complexity optimum is established. The original question's unspecified height and equivalence conventions remain material; the explicit unit is a concrete calculation, not a claimed complete answer.
\paragraph{Integration review.} Codex (GPT-6; exact serving revision not exposed) reviewed the written arguments and expanded layout and source-reference wrappers. Repository metadata and the production web renderer were checked using Python 3.13 and JavaScript with MathJax 3.2.2. This model-assisted review is not independent human proof verification; no PDF compilation or formal proof checking is asserted.

\section{Completion Estimate}
\noindent\textbf{COMPLETION: 0\%}
The calculations below the preserved source statement investigate a concrete route,
but establish no substantiated new progress toward the entire original target.
Known special cases and failed sufficient conditions are not credited as a solution.
\section{References}
The original source references and catalogue attribution are preserved in the
Sources section above. No additional source page was successfully retrieved for
this attempt; the retrieval limitation is stated in Literature and context check.
\section{Model and contributor disclosure}
OpenAI gpt-6-astra, requested reasoning configuration ultra, generated this research
attempt on 2026-10-11. The exact serving revision is not exposed. Codex GPT-6
performed orchestration, repository-format review, and written mathematical review at the declared scope; its exact serving revision
is also not exposed. The mathematical route, boundary checks and final adversarial
self-audit followed contribution prompts A then B in
\texttt{docs/contribution\_prompts.md}, after reading \texttt{CONTRIBUTING.md}.
Python 3.14.2 assembled the standalone documents; any mathematical computation is
specified in Verification. No theorem prover was used. No human mathematical
authorship, substantial human proof correction, or independent human proof checking
is asserted. Model-assisted review, source retrieval, and structural checks do not
establish correctness or novelty. Full private reasoning transcripts are unavailable
as reproducibility artifacts; the complete arguments, inputs and outcomes actually
used are stated here or in the companion batch audit.
\section{Final status}
Unresolved as a new research contribution. This is a concrete route investigation,
not a claim of new complete resolution or a current-status certification. A known
answer reproduced for context does not become a new contribution. The catalogue's
source-status label is preserved; this attempt status does not assert that a
source-labelled solved question is open.
\end{document}
